\section*{Chapitre \ref{ch:iv1:timed-test-3} --- Épreuve chronométrée 3}
\markboth{\protect\omTocMark{}{Réponses~: Épreuve chronométrée 3}}{\protect\omTocMark{}{Réponses~: Épreuve chronométrée 3}}

\begin{solution}{iq:iv1:timed-test-3:1}
La moyenne vaut $7.3$ boîtes.

Notons $T_i$ l'attente du jouet $i$. Le temps pour obtenir tous les jouets est $\max T_i$, et l'inclusion--exclusion\index{inclusion-exclusion} sur les minimums donne
\[
\mathbb{E}\max = \sum_i\frac{1}{p_i} - \sum_{i < j}\frac{1}{p_i + p_j} + \frac{1}{p_1 + p_2 + p_3}.
\]
La première somme vaut $2 + 3 + 6 = 11$. Les paires ont pour valeurs $\frac56$, $\frac23$ et $\frac12$, qui donnent $1.2 + 1.5 + 2 = 4.7$. Le dernier terme vaut 1. Donc la moyenne vaut $11 - 4.7 + 1 = 7.3$.

Vérification de bon sens~: elle doit dépasser 6, l'attente du jouet le plus rare, et c'est bien le cas, de peu. Avec des probabilités égales, on obtiendrait $5.5$, donc des probabilités inégales coûtent plus cher.

\iqwrong{Utiliser la formule à probabilités égales $3(1 + \frac12 + \frac13)$, qui donne $5.5$.}
\end{solution}

\begin{solution}{iq:iv1:timed-test-3:2}
Un total de 10 est plus probable~: $\frac{27}{216}$ contre $\frac{25}{216}$.

Les deux totaux admettent six partitions non ordonnées, ce qui explique qu'ils puissent sembler égaux. Mais les partitions n'ont pas le même nombre d'ordres possibles\index{dénombrement des issues ordonnées}.

Pour 9~: 126 (6), 135 (6), 144 (3), 225 (3), 234 (6), 333 (1), soit un total de 25.

Pour 10~: 136 (6), 145 (6), 226 (3), 235 (6), 244 (3), 334 (3), soit un total de 27.

Contrôle~: par la symétrie $x \to 7 - x$, les totaux 10 et 11 ont la même probabilité, de même que 9 et 12. Une énumération des 216 issues confirme les deux nombres.

\iqwrong{Compter les partitions comme équiprobables, ce qui traite 333 comme 126. Les issues sont les 216 triplets ordonnés, et seuls eux sont équiprobables.}
\end{solution}

\enlargethispage{-3\baselineskip}%
\begin{solution}{iq:iv1:timed-test-3:3}
Les deux valent $1.75$ bit.

La construction de Huffman\index{code de Huffman} fusionne les deux symboles les plus rares et recommence. Elle donne les longueurs de code 1, 2, 3 et 3. La longueur moyenne vaut $\frac12 \cdot 1 + \frac14 \cdot 2 + \frac18 \cdot 3 + \frac18 \cdot 3 = 1.75$.

L'entropie\index{entropie} vaut $\sum p\log_2(1/p)$, avec les mêmes nombres, puisque chaque $\log_2(1/p)$ est un entier. Le code atteint donc exactement la borne théorique.

Comparons avec un code à longueur fixe, qui prend 2 bits. Le gain vient de l'usage de mots courts pour les symboles fréquents.

\iqfollow{Dans quel cas le code n'atteint-il pas l'entropie~?}{Quand les probabilités ne sont pas des puissances de $\frac12$. Pour une source de probabilités $0.9$ et $0.1$, le code prend 1 bit par symbole, alors que l'entropie vaut $0.47$. Coder des blocs de symboles comble cet écart.}
\end{solution}

\begin{solution}{iq:iv1:timed-test-3:4}
La probabilité vaut $\frac12(1 - \ln 2) \approx 0.153$.

Conditionnons sur $X = x$\index{probabilités totales}. Pour $x \le \frac12$, l'événement est impossible. Pour $x > \frac12$, il a pour probabilité $(x - \frac12)/x = 1 - \frac{1}{2x}$. Donc
\[
P = \int_{1/2}^1\left(1 - \frac{1}{2x}\right)dx = \frac12 - \frac12\ln 2.
\]

Une seconde voie utilise la densité de $Y$, qui vaut $-\ln y$ sur $(0, 1)$, obtenue en intégrant $1/x$ de $y$ à 1. Alors $P = \int_{1/2}^1(-\ln y)\,dy = \frac12 - \frac12\ln 2$, le même résultat.

La moyenne de $Y$ vaut $\frac14$, donc une probabilité de $0.15$ de dépasser $\frac12$ est raisonnable.

\iqwrong{Utiliser la moyenne de $X$, $\frac12$, comme si $Y$ était uniforme sur $[0, \frac12]$, ce qui donne 0. Une valeur ne peut pas remplacer une loi ici.}
\end{solution}

\begin{solution}{iq:iv1:timed-test-3:5}
A gagne avec la probabilité $\frac38$.

À chaque tour complet, A gagne avec la probabilité $\frac16$. Si A échoue, B gagne avec la probabilité $\frac13$. La partie se poursuit donc à un tour supplémentaire avec la probabilité $\frac56 \cdot \frac23 = \frac59$\index{série géométrique}.

Notons $a$ la probabilité de A. Alors $a = \frac16 + \frac59a$, ce qui donne $a = \frac16 \cdot \frac94 = \frac38$.

Contrôle~: la probabilité de B vaut $\frac56 \cdot \frac13/(\frac49) = \frac58$, et les deux font 1. Commencer en premier ne suffit donc pas à compenser une cible deux fois plus facile. Avec deux cibles égales, un 6 chacune, A aurait $\frac{6}{11}$.

\iqwrong{Ne compter que le premier lancer de A, ou donner $\frac13$ en comparant directement les deux probabilités, $\frac16/(\frac16 + \frac13)$. Cela ignore qui lance en premier.}
\end{solution}

\begin{solution}{iq:iv1:timed-test-3:6}
Un triangle rectangle a la probabilité $\frac35$, et un triangle équilatéral la probabilité $\frac{1}{10}$.

Il y a $\binom63 = 20$ triangles. Les sommets de l'hexagone sont sur un cercle, donc un triangle a un angle droit exactement lorsque deux de ses sommets forment un diamètre\index{théorème de Thalès}. Il y a 3 diamètres. Chacun se combine avec l'un des 4 autres sommets, ce qui donne 12 triangles, tous distincts, puisqu'un triangle ne peut pas contenir deux diamètres.

Les triangles rectangles ont donc la probabilité $\frac{12}{20}$. Les triangles équilatéraux sont les deux formés par des sommets alternés, soit $\frac{2}{20}$. Les 6 restants sont isocèles à angle obtus~: trois sommets voisins. Une énumération confirme les effectifs.

\iqwrong{Chercher à calculer les angles de chaque cas. Le critère du diamètre donne le compte immédiatement.}
\end{solution}

\begin{solution}{iq:iv1:timed-test-3:7}
La probabilité vaut $\frac18$, et le nombre moyen de cycles vaut $H_8 = \frac{761}{280} \approx 2.72$.

Suivons le cycle à partir de 8. Il se referme au premier pas avec la probabilité $\frac18$, la valeur étant $\pi(8) = 8$. Sinon il se poursuit vers un autre élément, et se referme au pas suivant avec la probabilité $\frac17$, et ainsi de suite. La probabilité qu'il soit de longueur 3 vaut $\frac78 \cdot \frac67 \cdot \frac16 = \frac18$\index{longueur de cycle}. Toutes les longueurs de 1 à 8 sont équiprobables, par la même simplification.

Pour le nombre de cycles, utilisons que chaque élément $i$ d'un cycle de longueur $L_i$ contribue $1/L_i$, avec un total de 1 par cycle. Le compte vaut donc $\sum_i1/L_i$, et sa moyenne vaut $8 \cdot \frac18(1 + \frac12 + \dots + \frac18) = H_8$. Un dénombrement exhaustif de toutes les permutations donne le même résultat.

\iqfollow{Quelle est la longueur moyenne du cycle contenant un élément donné~?}{$(n + 1)/2 = 4.5$, la moyenne de la loi uniforme de 1 à 8. Remarquons qu'elle est supérieure à la longueur moyenne d'un cycle, $8/H_8 = 2.94$, un effet d'échantillonnage.}
\end{solution}

\begin{solution}{iq:iv1:timed-test-3:8}
Environ $0.196$ test par personne, soit une économie d'environ 80~\%.

Le lot coûte un test, soit $\frac{1}{10}$ par personne. Avec la probabilité $1 - 0.99^{10}$, le lot est positif, et chaque personne subit alors un test supplémentaire\index{tests de groupe}. La moyenne par personne vaut donc
\[
\frac{1}{10} + 1 - 0.99^{10} = 0.1 + 0.0956 = 0.1956.
\]
Pour la valeur, $0.99^{10} \approx e^{-0.1} \approx 0.904$.

La meilleure taille de lot équilibre les deux termes~: pour $f(k) = 1/k + 1 - 0.99^k$, le minimum se situe près de $k \approx 1/\sqrt{p} = 10$. Une recherche exacte montre que les lots de 10 et de 11 sont presque équivalents. Donc 10 est un bon choix.

\iqfollow{Si le taux monte à 30~\%, le recours aux lots est-il encore utile~?}{À peine. Un lot de 2 coûte $1.01$ par personne, le meilleur lot, de taille 3, coûte $0.99$, et les lots plus grands coûtent davantage. Il ne vaut donc plus la peine de s'en occuper.}
\end{solution}

\begin{solution}{iq:iv1:timed-test-3:9}
La probabilité vaut $\frac16(\frac76)^5 \approx 0.360$, et pour $n$ grand elle tend vers $\frac27$.

Notons $p_n$ la probabilité d'atteindre $n$, avec $p_0 = 1$. En conditionnant sur le dernier lancer\index{équation de renouvellement}, pour $n \le 6$,
\[
p_n = \frac16(p_{n - 1} + \dots + p_0).
\]
Donc $p_1 = \frac16$, et $p_{n + 1} = p_n + \frac16p_n = \frac76p_n$ pour $n \ge 1$ jusqu'à 5. Ainsi $p_6 = \frac16(\frac76)^5$. Une récurrence exacte en fractions donne ce nombre.

Pour $n$ grand, la théorie du renouvellement dit que les coups arrivent à un rythme d'un par pas moyen, et le pas moyen vaut $3.5$. Donc $p_n \to 1/3.5 = \frac27 = 0.286$.

Notons que $p_6$ est le plus grand des $p_n$~: à partir de 6, la borne sur le nombre de termes précédents dans la récurrence entre en jeu.

\iqwrong{Énumérer les compositions de 6 et pondérer chacune par $(\frac16)^k$. Cela fonctionne, mais c'est la forme lente de la même réponse.}
\end{solution}

\begin{solution}{iq:iv1:timed-test-3:10}
L'aire moyenne vaut $\frac16$, et l'aire dépasse $\frac{3}{16}$ avec une probabilité de $\frac12$.

Avec la cassure en $x$, uniforme sur $[0, 1]$, l'aire vaut $x(1 - x)$, dont la moyenne est $\frac12 - \frac13 = \frac16$. L'aire dépasse $\frac{3}{16}$ lorsque $(x - \frac14)(x - \frac34) < 0$, c'est-à-dire lorsque $x$ se trouve entre $\frac14$ et $\frac34$, et cela a une probabilité de $\frac12$.

Une vérification~: le morceau le plus court $S$ a pour moyenne $\frac14$ et le plus long $L$ a pour moyenne $\frac34$, avec un produit $\frac{3}{16}$. Or $L = 1 - S$, donc leur covariance\index{covariance} vaut $-\operatorname{Var}(S) = -\frac{1}{48}$, et $\frac{3}{16} - \frac{1}{48} = \frac16$.

\iqwrong{Donner $\frac{3}{16}$, le produit des morceaux moyens, ou $\frac14$, l'aire à la cassure moyenne. La première réponse ignore que les morceaux sont corrélés négativement, et la seconde est la plus grande aire possible, atteinte par un seul point de cassure.}

\iqwhy{La moyenne d'un produit n'est le produit des moyennes qu'en l'absence de corrélation. La moyenne du prix multiplié par le volume n'est pas la moyenne du prix multipliée par la moyenne du volume.}
\end{solution}

\begin{solution}{iq:iv1:timed-test-3:11}
Le straddle coûte environ $6.4$, et le mouvement quotidien de seuil de rentabilité est d'environ 1, soit 1~\%.

L'option d'achat vaut $0.4 \times 100 \times 0.16 \times 0.5 = 3.2$. À la monnaie, avec des taux nuls, le put et le call ont la même valeur, par parité. Donc le straddle coûte $6.4$.

Le facteur provient de $2\phi(0)\frac12 \approx 0.4$\index{approximation à la monnaie}~: le call vaut $S(2\Phi(\frac12\sigma\sqrt{T}) - 1)$, et il est linéaire en un petit argument.

Un portefeuille couvert en delta rapporte $\frac12\Gamma S^2(r^2 - \sigma^2\Delta t)$ par pas, de sorte qu'il atteint le seuil de rentabilité lorsque le rendement quotidien égale $\sigma\sqrt{\Delta t} = 16\%/\sqrt{252} = 1.01\%$\index{gamma}. C'est la règle du 16~: la volatilité annuelle divisée par 16 donne le mouvement quotidien.

\iqwrong{Comparer le seuil de rentabilité à l'échéance, $\pm 6.4$, avec les mouvements quotidiens. Le résultat d'un détenteur couvert dépend de la somme des carrés des mouvements, et non de l'endroit où l'action finit.}

\iqfollow{Combien coûte le straddle à un an~?}{$6.4 \times 2 = 12.8$, puisque l'échelle de temps double la racine.}
\end{solution}

\begin{solution}{iq:iv1:timed-test-3:12}
La moyenne vaut $\frac{11}{3}$.

Utilisons des indicatrices\index{variables indicatrices}. Une position intérieure, qui possède deux voisins, a la plus grande des trois valeurs aléatoires distinctes avec une probabilité de $\frac13$. Une position d'extrémité, qui possède un seul voisin, a une probabilité de $\frac12$.

Il y a 8 positions intérieures et 2 extrémités, donc la moyenne vaut $\frac83 + 1 = \frac{11}{3}$. La force brute pour de petits $n$ suit la formule générale $(n + 1)/3$.

Le résultat correspond à environ une position sur trois, soit le taux auquel une série aléatoire produit un pic. Ainsi, si un graphique semble présenter moins de points de retournement que cela, la série possède une tendance ou une mémoire.

\iqfollow{Quel est le nombre de séquences de hausses, pour $n = 10$~?}{Le nombre de hausses a pour moyenne $4.5$, et les séquences maximales de hausses sont en nombre supérieur d'une unité aux baisses, soit $5.5$ en moyenne. Ce sont les mêmes indicatrices, avec des paires à la place des triplets.}
\end{solution}
